% Imporditud: tools/import_eksperimendid.py
% Allikas: Eesti füüsikaolümpiaadi eksperimendi ülesannete
%   kogu (Rael Kalda, 2023), ülesanne Ü113, lahendus L113.
\setAuthor{Jaan Kalda}
\setRound{lõppvoor}
\setYear{2015}
\setNumber{G E1}
\setDifficulty{3}
\setTopic{dünaamika}
\setVahendid{kaks ühesugust nööri otsa kinnitatud elastset palli, statiiv, millimeeterpaber, kirjaklambrid}
\setPunktid{14}
\setAllikas{Eksperimendi kogu Ü113, lahendus lk 88}
\setLahendusseis{toores}

\prob{Kaks palli}
Määrata, mitu protsenti kineetilisest energiast muutub kahe palli omavahelisel põrkel soojuseks.

\hint

\solu
Riputame pallid paralleelsete pendlitena statiivi külge sarnaselt Newtoni hällile.

Mõõdame pendlite pikkused, olgu need L = 60 cm. Kergitame esimese kuuli üles

teatud kõrgusele, kusjuures külgsuunaline nihe olgu $x_0$ = 30 cm. Laseme selle lii-

kuma nii, et kuul lendab vastu teist (algselt paigalseisvat) kuuli. Kui põrge oleks

absoluutselt elastne, siis vahetaksid kuulid kiiruseid ehk esimene kuul jääks seis-

ma. Kuna põrkel läheb osa energiast kaduma, siis omandab see siiski teatud kii-

ruse, mille põhjal teemegi kindlaks põrkel kaotatud energia. Püüame teise (algselt

paigal olnud) kuuli, mis nüüd märkimisväärse kiiruse omandas, käega kinni ning

teeme kindlaks, millise amplituudiga jääb võnkuma esimene kuul. Olgu ampli-

tuud $x_1$ = 1,5 cm. Paneme tähele, et meetod eeldab täpselt tsentraalset põrget. Kui

põrge ei ole tsentraalne, omandab kuul külgsuunalise kiiruse, mille väärtus sõltub

esmajoones põrke mitte-tsentraalsusest (ja vähem mitte-elastsusest). Sestap tuleb

veenduda, et esimene kuul hakkaks võnkuma enam-vähem enda esialgse liiku-

mise suunas. Kui võnkumissuund on oluliselt pöördunud, tuleb katset korrata se-

ni, kuni õnnestub enam-vähem tsentraalne põrge. Olgu esimese kuuli kiirus enne

põrget $v_0$. Massikeskme süsteemis lähenevad mõlemad kuulid teineteisele kiiruse-

ga $v_0$ . Kui põrkel kaduma läinud energia suhe akv2l0ug(us1el$-$isdsee1kki$-$iinreukes)t.eiKdliisvisr1eu=estneve2s0rughi1atse$-$sevvk0o.snLakaa--,

siis pärast põrget on massikeskme süsteemis

boratoorses süsteemis on kiirus v = $v_0$ $-$$v_1$ =

me potentsiaalsete energiate abil: v = $h_1$/$h_0$, kus $h_0$ = L $-$ $L_2$ $-$ x20 $\approx$ 80 mm $v_0$

ja analoogiliselt $h_1$ $\approx$ = 1 $-$ 2 v $\approx$ 0,903, millest k $\approx$ 0,19 0,19 mm. Seega 1 $-$ k $v_0$

ehk põrkel läheb kaduma umbes 19 \% kineetilisest energiast.
\probend
